\section{Conclusion}
\label{sec:conclusion}
Under matched budgets and predeclared rules on three public collections, the
results were as follows.
\begin{itemize}
  \item Ranking peaks by correlation with a label-free GMM segmentation
        outperformed gradient attribution through a VAE. This was statistically
        supported and stable across seeds.
  \item Neighbourhood input to the encoder gave no consistent, spatially
        specific gain. Spatial information helped as a ranking statistic.
  \item A linear segmentation matched the VAE, and reconstruction did not track
        peak quality.
  \item msiPL reproduced. S3PL's GBM score did not, and its seeds switch
        sections between distinct solutions.
\end{itemize}
The contribution is a controlled evaluation rather than a new peak picker. It
bounds where spatial information helps learned MSI peak ranking, and it shows
how much conclusions depend on fair budgets, model-free references and seed
replication.

Future work should:
\begin{itemize}
  \item validate a label-free number of segments;
  \item test a predeclared combination of segmentation correlation and
        Moran's~$I$ on new data;
  \item evaluate graph- and ion-image-based spatial models;
  \item test transfer of fitted models between sections.
\end{itemize}
